
@@M0@@
où @@M1@@ est un polynôme universel en les @@M2@@, @@M3@@, polynôme qui a priori est à coefficients rationnels. En fait, on constate [RRR] que ces polynômes sont à coefficients \uline{entiers}. Cela soulève la question de la validité d'une formule de Riemann--Roch ``sans dénominateurs''. En fait, une telle formule pour @@M4@@ ne peut être cherchée en général dans @@M5@@ lui-même, vu que l'homomorphisme @@M6@@ n'est pas compatible lui-même avec les @@M7@@-filtrations modulo décalage par @@M8@@ (4.6), mais l'est seulement à torsion près; de sorte qu'il n'est pas possible de définir un homomorphisme @@M9@@ donnant l'homomorphisme naturel @@M10@@, et permettant de donner un sens à la formule (3.1) comme une formule dans @@M11@@, et non seulement dans @@M12@@. Mais lorsque @@M13@@ est supposé régulier, donc @@M14@@ (IV 2.5), on peut munir @@M15@@ de la \uline{filtration topologique} par la codimension des supports, héritée de la filtration correspondante de @@M16@@. Il est plausible que cette filtration est compatible avec la structure d'anneau, fait qu'on peut prouver tout au moins [RRR] lorsque @@M17@@ est lisse sur un corps @@M18@@, en utilisant la théorie de l'anneau de Chow [4]. Quand il en est ainsi, le gradué associé
@@M19@@
est un \uline{anneau gradué}, et supposant également @@M20@@ régulier, on a fait ce qu'il fallait pour avoir un homomorphisme naturel de degré @@M21@@ de groupes gradués
@@M22@@
Comme la @@M23@@-filtration de @@M24@@ est plus fine que la filtration topologique par la codimension qu'on vient de considérer (X 1.3.3), on trouve un homomorphisme canonique de groupes gradués (et même d'anneaux gradués)
@@M25@@
qui est d'ailleurs un isomorphisme mod torsion (Exp. VIII). Ceci permet, comme dans [RRR], de considérer une théorie des classes de Chern sur @@M26@@ à valeurs dans @@M27@@, et de même sur @@M28@@, et de donner ainsi un sens à la formule (3.1) comme une formule dans @@M29@@. Sous cette forme, cette formule est prouvée dans [RRR] lorsque @@M30@@ est de \uline{caractéristique nulle}, tout au moins si @@M31@@ est lisse sur un corps (hypothèse servant uniquement à assurer que la filtration topologique de @@M32@@ est compatible avec la structure d'anneau). L'hypothèse de caractéristique nulle est utilisée
essentielle%
